% Part: incompleteness
% Chapter: arithmetization-syntax
% Section: coding-symbols

\documentclass[../../../include/open-logic-section]{subfiles}

\begin{document}

\olfileid[fr]{inc}{art}{cod}
\olsection{Codage des symboles}

Le langage de base~$\Lang L$ de la logique du premier ordre emploie les
symboles
\[
\lfalse \quad \lnot \quad \lor \quad \land \quad \lif \quad \lforall
\quad \lexists \quad \eq \quad ( \quad ) \quad , \quad \ltrue \quad \liff
\]
ainsi que des ensembles !!{enumerable}s de variables et de
!!{constant}s, et des ensembles !!{enumerable}s de !!{function}s et
de !!{predicate}s d'arité quelconque. Nous pouvons attribuer un
\emph{code} à chacun de ces symboles de façon que chaque symbole
reçoive un nombre unique et que deux symboles distincts ne reçoivent
jamais le même nombre. Cela est possible puisque l'ensemble de tous
les symboles est !!{enumerable} : il existe donc !!a{bijection}
entre cet ensemble et celui des nombres naturels. Mais nous voulons
aussi pouvoir retrouver par calcul le symbole à partir de son code,
ainsi que certaines informations le concernant, par exemple l'arité
d'!!a{function}. Il existe bien des façons d'y parvenir. En voici
une qui utilise des fonctions récursives primitives. Rappelons que
$\tuple{n_0, \dots, n_k}$ est le nombre codant la suite de nombres
$n_0$, \dots, $n_k$.

\begin{defn}
Si $s$ est un symbole de~$\Lang L$, son \emph{code de symbole}~$\scode s$
est défini comme suit\footnote{Complément éditorial : la source omet $\ltrue$ et $\liff$ de la liste et du tableau, bien qu'ils figurent parmi les variantes primitives de sa syntaxe. Nous leur réservons les codes $\tuple{0,11}$ et $\tuple{0,12}$, sans modifier les autres codes. Les arités des familles de fonctions et de prédicats sont positives, comme dans la définition du langage standard ; nous rendons cette restriction explicite dans les tests de codage.} :
\begin{enumerate}
\item Si $s$ est l'un des symboles logiques, $\scode s$ est donné
  par le tableau suivant :
\[
\begin{array}{cccccccccc}
  \lfalse & \lnot & \lor & \land & \lif & \lforall \\
\tuple{0, 0} & \tuple{0, 1} & \tuple{0, 2} & \tuple{0, 3} &
\tuple{0, 4} & \tuple{0, 5} \\
\lexists & \eq & ( & ) & , & \ltrue & \liff\\
\tuple{0, 6} & \tuple{0, 7} &
\tuple{0, 8} & \tuple{0, 9} & \tuple{0, 10} & \tuple{0, 11} & \tuple{0, 12}
\end{array}
\]
\item Si $s$ est la $i$-ième variable $\Obj v_i$, alors
  $\scode s = \tuple{1, i}$.
\item Si $s$ est la $i$-ième !!{constant}~$\Obj c_i$, alors
  $\scode s = \tuple{2, i}$.
\item Si $s$ est le $i$-ième !!{function} d'arité~$n$, $\Obj f_i^n$,
  alors $\scode s = \tuple{3, n, i}$.
\item Si $s$ est le $i$-ième !!{predicate} d'arité~$n$, $\Obj P_i^n$,
  alors $\scode s = \tuple{4, n, i}$.
\end{enumerate}
\end{defn}

\begin{prop}
Pour $n>0$, les relations suivantes sont récursives primitives :
\begin{enumerate}
\item $\fn{Fn}(x, n)$ si et seulement si $x$ est le code de
  $\Obj f^n_i$ pour un certain~$i$, c'est-à-dire que $x$ est le code
  d'!!a{function} d'arité~$n$.
\item $\fn{Pred}(x, n)$ si et seulement si $x$ est le code de
  $\Obj P^n_i$ pour un certain~$i$, ou bien si $x$ est le code de $\eq$
  et $n = 2$, c'est-à-dire que $x$ est le code d'!!a{predicate}
  d'arité~$n$.
\end{enumerate}
\end{prop}

\begin{defn}
Si $s_0, \dots, s_{n-1}$ est une suite de symboles, son
\emph{numéro de G\"odel} est
$\tuple{\scode{s_0}, \dots, \scode{s_{n-1}}}$.
\end{defn}

\begin{explain}
Il faut distinguer les \emph{codes} des \emph{numéros de G\"odel}.
Par exemple, la variable~$\Obj v_5$ a pour code~$\scode{\Obj
  v_5} = \tuple{1, 5} = 2^2\cdot 3^6$. Mais, considérée comme un
terme, la variable~$\Obj v_5$ est aussi une suite de symboles, de
longueur~$1$. Le \emph{numéro de G\"odel}~$\Gn{\Obj v_5}$ du
\emph{terme}~$\Obj v_5$ est $\tuple{\scode{\Obj v_5}} = 2^{\scode{\Obj
    v_5} + 1} = 2^{2^2\cdot 3^6 + 1}$.
\end{explain}

\begin{ex}
Rappelons que, si $k_0$, \dots, $k_{n-1}$ est une suite de nombres,
le code de la suite $\tuple{k_0, \dots, k_{n-1}}$, dans le codage
par puissances de nombres premiers, est
\[
2^{k_0+1}\cdot3^{k_1+1}\cdot \dots \cdot p_{n-1}^{k_{n-1}+1},
\]
où $p_i$ est le $i$-ième nombre premier, avec $p_0 = 2$. Ainsi, par
exemple, la formule $\eq[\Obj v_0][\Obj 0]$, ou, plus explicitement,
${\eq}(\Obj v_0,\Obj c_0)$, a pour numéro de G\"odel
\[
\tuple{\scode{\eq},\scode{(},\scode{\Obj v_0},\scode{,}, \scode{\Obj
    c_0},\scode{)}}.
\]
Ici, $\scode{\eq}$ vaut $\tuple{0,7} = 2^{0+1}\cdot
3^{7+1}$, $\scode{\Obj v_0}$ vaut $\tuple{1,0} = 2^{1+1}\cdot3^{0+1}$,
et ainsi de suite. Par conséquent, $\Gn{=(\Obj v_0,\Obj c_0)}$ vaut
\begin{multline*}
2^{\scode{=} + 1}\cdot 3^{\scode{(}+1}\cdot 5^{\scode{\Obj v_0}+1}
\cdot 7^{\scode{,} + 1} \cdot 11^{\scode{\Obj c_0}+1} \cdot
13^{\scode{)}+1} = \\
2^{2^1\cdot 3^8 + 1}\cdot 3^{2^1\cdot 3^9+1}\cdot 5^{2^2\cdot 3^1+1}
\cdot 7^{2^1\cdot 3^{11} + 1} \cdot 11^{2^3\cdot3^1+1} \cdot
13^{2^1\cdot3^{10}+1} = \\
2^{13\,123}\cdot 3^{39\,367}\cdot 5^{13}\cdot 7^{354\,295}\cdot11^{25}\cdot13^{118\,099}.
\end{multline*}
\end{ex}

\end{document}
